% !TeX root = ../main.tex
% Add the above to each chapter to make compiling the PDF easier in some editors.

\chapter{Algorithms and Listings}

\begin{center}
    \begin{lstlisting}[
        language=python,
        caption={Pseudocode for the overview and initialization phase of the greedy strategy.},
        label={fig:greedy_overview_pseudocode},
        numbers=left,
    ]
    def FIND_OPTIMAL_PERMUTATION(network, contraction_path, k):
        persistent_path = build_persistent_contraction_path(network, 
            contraction_path)
        initial_perms = identity_permutations(network.tensors)
        intermediate_perms = identity_permutations(
            persistent_path.intermediate_results)

        if persistent_path.num_steps == 0:
            return materialize_operations(initial_perms, 
                intermediate_perms, contraction_path)

        tree, leaf_to_node, step_to_node = build_tree_maps(persistent_path)
        peak_nodes = find_peak_nodes(network, 
            persistent_path, leaf_to_node, step_to_node, k)

        desired_layouts = {}
        for node in peak_nodes:
            desired_layouts[node.id] = compute_peak_target_layout(node, 
                desired_layouts)

        plan_recursive_layouts(tree.root, desired_layouts, 
            frozen_nodes=peak_nodes)
        return materialize_operations(initial_perms, intermediate_perms, 
            contraction_path)
    \end{lstlisting}
\end{center}

\newpage

\begin{center}
    \begin{lstlisting}[
        language=python,    
        caption={Pseudocode for selecting the top-$k$ peak tensors from the contraction tree.},
        label={fig:greedy_peak_nodes_pseudocode},
        numbers=left,
    ]
    def FIND_PEAK_NODES(network, persistent_path, 
        leaf_to_node, step_to_node, k):
            candidates = []
            candidates += largest_tensors(network)
            candidates += largest_intermediate_tensors(persistent_path)

            sort_descending(candidates, key=memory_size)
            selected = candidates.take_first(k)

            return sort_by_leaf_status(selected)
    \end{lstlisting}
\end{center}

\begin{center}
    \begin{lstlisting}[
        language=python,
        caption={Pseudocode for computing the target layout of a peak tensor.},
        label={fig:greedy_peak_layout_pseudocode},
        numbers=left,]
    def COMPUTE_PEAK_TARGET_LAYOUT(node, desired_layouts):
        if node.is_leaf:
            return node.input_indices

        if node.is_root:
            return sort_by_size(node.input_indices)

        parent = node.parent
        left, right, _ = get_step_tensors(parent.contraction_step)
        contracted = get_contracted_indices(node.tensor, sibling(node))
        free = node.input_indices - contracted

        if desired_layouts.contains(parent.id):
            free_order = sort_by_layout(free, desired_layouts[parent.id])
        else:
            free_order = sort_by_size(free)

        if desired_layouts.contains(sibling(node).id):
            contracted_order = sort_by_layout(contracted, 
                desired_layouts[sibling(node).id])
        else:
            contracted_order = sort_by_size(contracted)

        if node is left child:
            return concatenate(free_order, contracted_order)
        else:
            return concatenate(contracted_order, free_order)
    \end{lstlisting}
\end{center}

\begin{center}
    \begin{lstlisting}[
        language=python,
        caption={Pseudocode for checking whether the free-index order is representable without slicing.},
        label={fig:greedy_order_check_pseudocode},
        numbers=left,
    ]
    def HAS_LEFT_THEN_RIGHT_ORDER(desired_free, left_free, right_free):
        seen_right = false
        for index in desired_free:
            if index in right_free:
                seen_right = true
            else if index in left_free and seen_right:
                return false
        return true
    \end{lstlisting}
\end{center}

\begin{center}
    \begin{lstlisting}[
        language=python,
        caption={Pseudocode for selecting sliced indices when free-index order is interleaved.},
        label={fig:greedy_select_sliced_pseudocode},
        numbers=left,
    ]
    def SELECT_SLICED_INDICES(desired_free, left_free, right_free):
        best_keep = -1
        best_split = 0
        prefix_left = compute_prefix_counts(desired_free, left_free)
        suffix_right = compute_suffix_counts(desired_free, right_free)

        for split from 0 to length(desired_free):
            keep = prefix_left[split] + suffix_right[split]
            if keep > best_keep:
                best_keep = keep
                best_split = split

        sliced = {}
        for i, index in enumerate(desired_free):
            keep_left = (i < best_split and index in left_free)
            keep_right = (i >= best_split and index in right_free)
            if not (keep_left or keep_right):
                sliced.add(index)

        return sliced
    \end{lstlisting}
\end{center}

\begin{center}
    \begin{lstlisting}[
        language=python,
        caption={Pseudocode for the recursive layout planning phase.},
        label={fig:greedy_recursive_planning_pseudocode},
        numbers=left,
    ]
    def PLAN_NODE(node, frozen_nodes, desired_layouts): 
        if node.is_leaf:
            return

        left, right, result = get_step_tensors(node.step)
        contracted = get_contracted_indices(left, right)
        left_free = left.indices - contracted
        right_free = right.indices - contracted

        desired_free = determine_desired_free_order(node, desired_layouts)
        if HAS_LEFT_THEN_RIGHT_ORDER(desired_free, left_free, right_free):
            sliced = {}
        else:
            sliced = SELECT_SLICED_INDICES(desired_free, left_free, right_free)

        left_target = build_left_target(left_free, 
            sliced, contracted, desired_free)
        right_target = build_right_target(right_free, 
            sliced, contracted, desired_free)
        result_target = build_result_target(left_free, 
            right_free, sliced, desired_free)

        if node.left not in frozen_nodes:
            assign_layout(node.left, left_target)
        if node.right not in frozen_nodes:
            assign_layout(node.right, right_target)
        if node not in frozen_nodes:
            assign_layout(node, result_target)

        PLAN_NODE(node.left, frozen_nodes, desired_layouts)
        PLAN_NODE(node.right, frozen_nodes, desired_layouts)
    \end{lstlisting}
\end{center}

\newpage
\begin{center}
    \begin{lstlisting}[
        language=python,
        caption={Pseudocode for assigning a permutation to a tensor node after layout planning.},
        label={fig:greedy_assign_permutation_pseudocode},
        numbers=left,
    ]
    def ASSIGN_PERMUTATION(node, target_layout):
        tensor = get_node_tensor(node)
        permutation = compute_permutation(tensor.indices, target_layout)

        if node.is_leaf:
            initial_permutations[node.position] = permutation
        else:
            intermediate_permutations[node.step] = permutation

        desired_layouts[node.id] = target_layout
    \end{lstlisting}
\end{center}